\chapter{Инцидентный перенос, марковский вес и индексный quotient}
\label{ch:v7-incidence-transfer-markov-weight-gate}

\section{Что означала лишняя отрицательная мода}

Предыдущий гейт обнаружил, что честный source-corner даёт сигнатуру
$(8,0,19)$ вместо требуемой $(7,0,20)$. Это не следует читать как почти
правильный восьмой корень. Явная диагонализация показывает, что дополнительная
отрицательная мода лежит на нежелательном ребре $L_LX_R$:
\begin{equation}
 \lambda_{\rm extra}=-\frac25,
 \qquad
 v_{\rm extra}=\operatorname{Im}(L_LX_R)_{w=1}.
 \label{eq:v7-transfer-extra-source-mode}
\end{equation}
Причина структурна. Кривизна $A^*A$ помнит общий источник стрелок, но
суммирует по их целям. Поэтому source-only чтение теряет часть данных о том,
в какой правый блок приходит стрелка, и запускает запрещённого конкурента.

Это совпадает с уроком Тома V: диагональные условные ожидания читают
состояния, тогда как взаимодействие хранится во внедиагональном углу
связывающей алгебры. Следовательно, исправлять $(8,0,19)$ новым скалярным
весом нельзя; нужно проверить двусторонний перенос.

\section{Полярная часть существующей инцидентности}

Однопоколенческий фоновый оператор имеет тип
\begin{equation}
 A_0:\mathbb C^{11}\longrightarrow\mathbb C^{10},
 \qquad \rank A_0=10.
 \label{eq:v7-transfer-reference-rank}
\end{equation}
Его полярная часть определяется без нового коэффициента:
\begin{equation}
 U=A_0(A_0^*A_0)^{-1/2}_{\operatorname{supp}},
 \qquad
 UU^*=I_{10},
 \qquad
 P=U^*U,
 \qquad Q=I_{11}-P.
 \label{eq:v7-transfer-polar-coisometry}
\end{equation}
Здесь
\begin{equation}
 \rank P=10,
 \qquad \rank Q=1,
 \qquad
 21=10_{\rm source\ support}+10_{\rm target}+1_{\rm defect}.
 \label{eq:v7-transfer-index-decomposition}
\end{equation}
Именно это разложение, а не число $11/21$ само по себе, является операторным
содержанием размерностной зацепки.

Для переменного прямоугольного поля $A$ введём две относительные
Gram-кривизны
\begin{equation}
 C_s(A)=A^*A-A_0^*A_0,
 \qquad
 C_t(A)=AA^*-A_0A_0^*.
 \label{eq:v7-transfer-two-gram-curvatures}
\end{equation}

\section{Однократное quotient-чтение}

Полярная коизометрия задаёт унитальную вполне положительную карту
\begin{equation}
 T_U(X,Y)=\frac12\left(UXU^*+Y\right):
 M_{11}(\mathbb C)\oplus M_{10}(\mathbb C)\longrightarrow M_{10}(\mathbb C).
 \label{eq:v7-transfer-ucp-channel}
\end{equation}
Она является средним двух вполне положительных отображений и удовлетворяет
$T_U(I_{11},I_{10})=I_{10}$.

Рассмотрим однократно прочитанную кривизностную норму
\begin{equation}
 \mathcal S_{\rm quot}(A)
 =\frac12\left\|T_U(C_s(A),C_t(A))\right\|_{\rm HS}^2.
 \label{eq:v7-transfer-quotient-action}
\end{equation}
В нуле её гессиан удовлетворяет точному тождеству
\begin{equation}
 \Hess_0\mathcal S_{\rm quot}
 =\frac12\Hess_0\mathcal S_V.
 \label{eq:v7-transfer-half-hessian-identity}
\end{equation}
Таким образом, ранее условный коэффициент $\beta=1/2$ возникает локально не
из умножения на число, а из усреднения двух полярно отождествлённых Gram-
концов и однократного чтения их общего образа.

После сложения с уже выведенной рёберной Hodge-нормой аудит даёт
\begin{align}
 \operatorname{sig}\Hess_0\mathcal S_{\rm tot}^{\rm quot}
  &=(7,0,20),
 &\lambda_{\rm heavy}^{\min}&=\frac25,
 \label{eq:v7-transfer-quotient-origin-pass}\\
 \operatorname{sig}\Hess_{A_0}\mathcal S_{\rm tot}^{\rm quot}
  &=(0,0,27),
 &\lambda_{\rm vac}^{\min}&=4.5827630148\ldots.
 \label{eq:v7-transfer-quotient-vacuum-pass}
\end{align}
Это первый операторный локальный проход без вставленного секторного числа.

\section{Почему полное условное ожидание снова проваливается}

Карта~\eqref{eq:v7-transfer-ucp-channel} сама по себе не является
сохраняющим полный след условным ожиданием на исходной алгебре. Настоящее
ожидание на подалгебру отождествлённых углов имеет вид
\begin{equation}
 E_U(X\oplus Y)
 =\left(U^*ZU+\Tr(QXQ)Q\right)\oplus Z,
 \qquad
 Z=T_U(X,Y).
 \label{eq:v7-transfer-trace-preserving-expectation}
\end{equation}
Оно унитально, вполне положительно, идемпотентно и сохраняет полный след на
$M_{11}\oplus M_{10}$. Но общий десятимерный образ $Z$ появляется в нём в
обоих согласованных углах. Ранга-один дефект не даёт квадратичного вклада в
нуле, поэтому
\begin{equation}
 \Hess_0\frac12\|E_U(C_s\oplus C_t)\|_{\rm HS}^2
 =2\Hess_0\mathcal S_{\rm quot}.
 \label{eq:v7-transfer-expectation-double-count}
\end{equation}
Полный сохраняющий след вариант возвращает эффективный вес $\beta=1$ и
сигнатуру
\begin{equation}
 (n_-,n_0,n_+)_{E_U}=(21,0,6).
 \label{eq:v7-transfer-expectation-no-go}
\end{equation}

\begin{theorem}[Локальный проход полярного quotient и no-go полного ожидания]
Полярная часть уже существующего фонового оператора канонически разбивает
физический носитель как $10+10+1$. Однократное унитальное вполне
положительное чтение двух согласованных десятимерных углов выводит локальный
вес $1/2$, правильную сигнатуру нуля и положительный гессиан целевого
вакуума. Полное сохраняющее след условное ожидание считает общий угол в двух
диагональных копиях, возвращает вес $1$ и проваливает селекцию.
\end{theorem}

\begin{nogocriterion}[Запрет преждевременного объявления общего действия]
Нельзя отождествлять однократный quotient-readout с полным родительским
следом, пока не доказано, что физический фактор форм, Real-условие или
BRST--BV-редукция действительно оставляют одну согласованную копию
$M_{10}$ и отдельно сохраняют ранга-один индексный дефект. Без такой
теоремы карта $T_U$ является положительным локальным кандидатом, а не
завершённым действием.
\end{nogocriterion}

\begin{protocol}
\begin{enumerate}
 \item ранг фоновой инцидентности: \textbf{$10$ в типе $10\times11$};
 \item полярная коизометрия и проекторы $P,Q$: \textbf{выведены};
 \item истинное разложение размерности: \textbf{$10+10+1$};
 \item лишняя source-мода: \textbf{нежелательное ребро $L_LX_R$};
 \item однократная UCP-карта: \textbf{канонична относительно $A_0$};
 \item тождество гессианов: \textbf{$H_{\rm quot}=H_V/2$ в нуле};
 \item сигнатура нуля quotient-действия: \textbf{$(7,0,20)$};
 \item устойчивость целевого вакуума: \textbf{получена локально};
 \item полный сохраняющий след expectation: \textbf{$(21,0,6)$, провал};
 \item физическое происхождение редуцированного quotient: \textbf{открыто};
 \item следующий гейт: \textbf{индексно-дефектный редуцированный
 связывающий след}.
\end{enumerate}
\end{protocol}

Машинный сертификат сохранён в
\path{s2t_v7_incidence_transfer_markov_weight_gate_results.json}.

\begin{thebibliography}{9}
\bibitem{v7-transfer-kodaka-teruya}
K.~Kodaka, T.~Teruya,
\emph{The Strong Morita Equivalence for Inclusions of $C^*$-Algebras and
Conditional Expectations for Equivalence Bimodules},
arXiv:1609.08263.
\bibitem{v7-transfer-knill}
O.~Knill,
\emph{The McKean--Singer Formula in Graph Theory},
arXiv:1301.1408.
\end{thebibliography}